% БҮЛЭГ 7 — ЁС ЗҮЙН АСУУДАЛ

\chapter{Ёс зүйн асуудал}
\label{Chapter7}

Энэ бүлэгт өгөгдлийн ёс зүй, annotation, bias, нууцлал, хос хэрэглээний эрсдэл, AI мэдүүлэлт гэсэн зургаан хэмжээсийг авч үзэж Ш-5 шаардлагыг хангана.

%-------------------------------------------------------------------------------
\section{Өгөгдлийн ёс зүй ба зөвшөөрөл}

Найман нийтэд нээлттэй эх сурвалжаас нийтийн сэтгэгдлүүдийг ашигласан тул хувь хүний зөвшөөрөл шаардахгүй. Хэрэглэгчийн нэр, профайл, холбоос зэрэг хувийн мэдээллийг (PII) автоматаар устгаж, ``Хувь хүний нууцын тухай хууль'' (2021) болон GDPR зарчимд нийцүүлэн зөвхөн судалгааны зорилгоор ашигласан. Өгөгдөл нь нийтэд нээлттэй бөгөөд үл таних болгосон (de-identified) тул Ёс зүйн хорооны (IRB) зөвшөөрөл шаардахгүй.

\section{Тэмдэглэгээ ба тэмдэглэгээ хийгчдийн хоорондын нийцлийн шинжилгээ}

Нийт 3 тэмдэглэгээ хийгч эхлэлийн 200 жишээг хамтад нь шошголж удирдамж боловсруулсан; гол багцын $\geq\!20\%$-г давхар шошголсон бөгөөд Флейссийн каппа Fleiss' $\kappa{=}0.72$ (өндөр буюу \textit{substantial} нийцэл) гарсан нь өгөгдлийн багцын найдвартай байдлыг баталсан. Санал зөрөлдөөнтэй тохиолдлуудыг гуравдагч тэмдэглэгээ хийгч шийдвэрлэсэн.

\section{Хазайлт (Bias) ба тэгш байдлын шинжилгээ}


\subsection{Ангиллын түвшний хазайлт}
\label{sec:class-bias}

Stage~1 загвар нь Саармаг ангиллын дээж хамгийн олон учраас тухайн ангилалд хэт дасах (overfit) болох, Stage~2 загвар Бүтээлч шүүмжлэлийн цөөн дээжийг дутуу таних эрсдэл бий. SMOTE болон Focal Loss-ын хослол нь энэ тэнцвэргүй байдлыг зохицуулж, цөөнх (minority) ангиллын сэргээлтийг (recall) нэмэгдүүлсэн.

\subsection{Демографийн хазайлт}
\label{sec:demographic-bias}

Найман эх сурвалжийн өгөгдөл нь нийгмийн бүлгүүдээр жигд бус тархалттай байх магадлалтай тул загвар тодорхой бүлэгт шударга бус хандах эрсдэлтэй; ирээдүйд илүү тэнцвэртэй өгөгдөл шаардлагатай.

\subsection{Бүтээлч шүүмжлэлийг Хортой сөрөг гэж буруу ангилах эрсдэл}
\label{sec:false-positive-risk}

Хамгийн нухацтай алдааны төрөл нь хуулийн хүрээнд хамгаалагдсан бүтээлч шүүмжлэлийг хортой гэж буруу ангилах бөгөөд энэ нь хэрэглэгчийн үзэл бодлоо илэрхийлэх эрхийг зөрчихөд хүргэж болзошгүй. Ангиллын матрицийн диагоналийн бус утгуудыг хянаж, алдаатай ангиллыг хүний хяналтаар шүүх хүний оролцоотой (human-in-the-loop) тогтолцоо шаардлагатай.

Тэнцвэргүй байдлыг засахын тулд SMOTE (суурь загварт), Focal Loss + жин + Sampler (MN-BERT-д) хэрэглэж, macro F1-ийг үндсэн метрик болгон ашигласан. Олон төрлийн эх сурвалжийг хослуулан демографийн хазайлтыг бууруулж, хуурамч эерэг (false positive) алдааны эсрэг хүний оролцоотой (human-in-the-loop) хяналтыг санал болгов.

\section{Нууцлал ба хэрэглэгчийн өгөгдлийн хамгаалалт}

Хэрэглэгчийн нууцлалыг хамгаалах дараах арга хэмжээг баримталсан:

\begin{enumerate}
  \item \textbf{Хувийн мэдээллийг (PII) хасах:} хэрэглэгчийн нэр, профайл, холбоосыг автоматаар устгаж зөвхөн анонимчлагдсан текст болон огноог хадгалав.
  \item \textbf{Хадгалалт:} түүхий өгөгдлийг локал орчинд хадгалан нийтэд дэлгэхгүй бөгөөд судалгаа дуусахад устгана.
  \item \textbf{Загварын нууцлал:} сургах үеийн давталтын тоог хязгаарлах, dropout ашиглах; алдааны шинжилгээнд бодит бус зохиомол өгүүлбэрүүд ашигласан.
\end{enumerate}

%-------------------------------------------------------------------------------
\section{Хос хэрэглээний эрсдэл (Dual-use risk)}

NLP систем нь хортой контентыг шүүх эерэг талтай боловч хууль ёсны үзэл бодлыг дарангуйлах сөрөг үр дагавартай байж болно. Ялангуяа улс төрийн шүүмжлэлийг буруу ангилах эрсдэлтэй. Эрсдэлийг бууруулахын тулд: ``зөвлөх'' үүрэг бүхий хүний оролцоотой систем (human-in-the-loop), тайлбарлах боломжтой байдал (explainability) болон тогтмол хүний хяналт-шалгалт (human audit) хийхийг зөвлөж байна.

\section{AI хэрэгслийн ашиглалт}
\label{sec:ai-disclosure}

Зааварчилгааны (Guideline) зүйл 7-д заасны дагуу: \textbf{Anthropic Claude} (алдаа засах, LaTeX форматлах) болон \textbf{Google Gemini} (санаа харьцуулах) зэргийг зөвхөн туслах үүрэгтэйгээр ашигласан; бүх тэмдэглэгээ, сургалт, шинжилгээний ажлыг зохиогч өөрийн локал орчинд бие даан гүйцэтгэсэн.

%-------------------------------------------------------------------------------
\section{Бүлгийн дүгнэлт}

Энэ бүлэгт зургаан хэмжээст ёс зүйн шинжилгээ (өгөгдлийн ёс зүй, тэмдэглэгээ, хазайлт, нууцлал, хос хэрэглээний эрсдэл, AI мэдүүлэлт) хийж, Ш-5 шаардлагыг хангав. Гол дүгнэлт: өгөгдлийг үл таних болгосон (de-identified) тул хууль, ёс зүйн зөвшөөрөл баталгаатай; $\kappa{=}0.72$ үзүүлэлт нь өгөгдлийн багц найдвартай болохыг илтгэж байна; Бүтээлч шүүмжлэлийг Хортой сөрөг гэж буруу ангилах эрсдэл байгаа тул хүний оролцоотой (human-in-the-loop) хяналтын тогтолцоо шаардлагатай.
